%%%%%%%%%%%%%%%%%%%%%%%%%%%%%%%%%%%%%%%%%%%%%%%%%%%%%%%%%%%%%%%%%%%%%%%%%%%%%
%% 03_arch_memory_concept.tex
%% Section: Memory Concept
%% Chapter: Architecture Concepts
%%
%% Variant A (RV_NoPipeline) only: chip variant summary and the
%% no-cache Harvard memory subsystem. Variant B has its own document.
%%
%% Integration in 00_rwurv64i_all.tex:
%%   \section{Memory Concept}
%%   \input{\TEXTarch 03_arch_history_tab07}
%%   \input{\TEXTarch 03_arch_memory_concept}
%%
%% Requires packages: tabularx, booktabs, tikz, float, xcolor, amsmath
%%%%%%%%%%%%%%%%%%%%%%%%%%%%%%%%%%%%%%%%%%%%%%%%%%%%%%%%%%%%%%%%%%%%%%%%%%%%%

%%%%%%%%%%%%%%%%%%%%%%%%%%%%%%%%%%%%%%%%%%%%%%%%%%%%%%%%%%%%%%%%%%%%%%%%%%%%%
\subsection{Chip Variants}
\label{subsec:chipVariants}
%%%%%%%%%%%%%%%%%%%%%%%%%%%%%%%%%%%%%%%%%%%%%%%%%%%%%%%%%%%%%%%%%%%%%%%%%%%%%

The \aschip project comprises two chip variants that implement the same
RV64I ISA but differ in CPU core and memory subsystem. This document covers
Variant~A only; Table~\ref{tab:asChipVariants} summarises it.

\begin{table}[H]
\caption{Variant A -- Summary}
\label{tab:asChipVariants}
\centering
\begin{tabularx}{\textwidth}{|l |X|}
  \hline
   & \textbf{Variant~A -- No-Cache} \\
  \hline
  \hline
  RTL top & \texttt{RV\_NoPipeline} \\
  \hline
  CPU core & \texttt{as\_cpux}: multi-cycle FSM, no pipeline (Section~\ref{sec:archcpu}) \\
  \hline
  Instruction memory & On-chip I-Mem (32\,KB), loaded via JTAG scan chain \\
  \hline
  Data memory & On-chip D-Mem (8\,KB), connected via Wishbone \\
  \hline
  External memory & None (no Flash, no QSPI controller) \\
  \hline
  Cache subsystem & None \\
  \hline
  Peripheral bus & Wishbone (D-Mem, GPIO0, CGU, UART0, ICU0) \\
  \hline
  Technology & X-FAB XH018, 180\,nm (Section~\ref{sec:archtech}) \\
  \hline
  Design goal & Minimal area; standalone execution from on-chip memory \\
  \hline
\end{tabularx}
\end{table}

%%%%%%%%%%%%%%%%%%%%%%%%%%%%%%%%%%%%%%%%%%%%%%%%%%%%%%%%%%%%%%%%%%%%%%%%%%%%%
\subsection{Variant A -- No-Cache Memory Subsystem}
\label{subsec:noCacheVariant}
%%%%%%%%%%%%%%%%%%%%%%%%%%%%%%%%%%%%%%%%%%%%%%%%%%%%%%%%%%%%%%%%%%%%%%%%%%%%%

Variant~A implements a simple Harvard memory subsystem with two on-chip SRAM
blocks and no cache.

\begin{itemize}
  \item \textbf{Instruction Memory (I-Mem, 32\,KB):}
    A single-port synchronous SRAM containing 8\,192 RISC-V 32-bit
    instructions (word-addressed).
    The CPU accesses I-Mem via a legacy flat bus.
    At reset, I-Mem is empty; the content is loaded via the JTAG scan chain
    (a dedicated instruction-memory shift register in the TAP).
    I-Mem is read-only during normal CPU execution.

  \item \textbf{Data Memory (D-Mem, 8\,KB):}
    A single-port synchronous SRAM with a Wishbone~B3 slave interface.
    Supports byte, halfword, word, and doubleword accesses controlled by the
    Wishbone \texttt{SEL} byte-enable signal.
    Holds \texttt{.data}, \texttt{.bss}, stack, and heap.
    Address range: \texttt{0x0000\_0000}--\texttt{0x0000\_1FFF}
    (8\,KB, mapped at base address 0).

  \item \textbf{Peripheral Bus:}
    All CPU data accesses are routed through a Wishbone master BPI.
    Address decoding selects D-Mem, GPIO0, CGU, UART0 or ICU0 by chip
    select (\texttt{csx\_s}).
\end{itemize}

There is no cache miss penalty in Variant~A: every fetch or load/store
completes in a fixed number of Wishbone cycles.
The trade-off is that all code and data must fit in the two 32\,KB + 8\,KB
on-chip memories. Both memories are XH018 SPRAM macros in silicon
(Section~\ref{subsec:sramSelectionA}).
